\NativeTitle{5}{第一单元素养评价（B）}{（四则混合运算）}{40}{5}{一 & 二 & 三 & 四 & 五}{ &  &  &  & }
\NativeSlot{285}{39}{%
\ExamGroup{一、填空。{\normalfont\kaishu\fontsize{10.8bp}{14bp}\selectfont （第3小题5分，其余各题每空1分，共22分）}}

% p005_s01

\ExamQuestion{1.}{5}%
计算\(12\times8+15\div3\)时，可以同时先算\mbox{（\hspace{1.8em}）}法和\mbox{（\hspace{1.8em}）}法，再算\mbox{（\hspace{1.8em}）}法，结果是\mbox{（\hspace{1.8em}）}。\par
}
\NativeSlot{358}{60}{%
% p005_s02

\ExamQuestion{2.}{5}%
按要求在算式\(32\times400-200\div25\)中添上括号，改变运算顺序。（不改变运算符号）\par\noindent （1）先减再除最后乘：\rule{10em}{.4pt}。\par

% p005_s03

\ExamQuestion{（2）}{5}%
先减再乘最后除：\rule{10em}{.4pt}。\par
}
\NativeSlot{459}{73}{%
% p005_s04

\ExamQuestion{3.}{7}%
烟台海产品特别有名。扇贝每千克12元，梭子蟹每千克96元，即食海鲜礼盒每盒的价格比扇贝和梭子蟹单价和的3倍多10元。即食海鲜礼盒每盒\mbox{（\hspace{1.8em}）}元。根据题意进行分析，将流程图填写完整。\par
\BookFigure{figures/p005_s04-01.tex}
}
\NativeSlot{582}{36}{%
% p005_s05

\ExamQuestion{4.}{5}%
630减去78的差乘14再除以7，商是\mbox{（\hspace{1.8em}）}。\par
}
\NativeSlot{609}{67}{%
% p005_s06

\ExamQuestion{5.}{5}%
对比算式，在圆圈里填上“\(>\)”“\(<\)”或“\(=\)”：\par\noindent ①\enspace \(150-75\div5\times3\ \bigcirc\ (150-75)\div5\times3\)。\par

% p005_s07

\noindent ②\enspace \(240\div(6+2)\times4\ \bigcirc\ 240\div[(6+2)\times4]\)。对比两组算式，我发现：括号能\BookBlank{}，从而改变算式的计算结果。\par
}
\NativeSlot{722}{42}{%
% p005_s08

\ExamQuestion{6.}{5}%
根据分步算式，写出综合算式：\par\noindent ①\enspace \(36\div4=9\)，\allowbreak{}\(12+9=21\)，\allowbreak{}\(21\times5=105\)。综合算式：\rule{9em}{.4pt}。\par

% p005_s09

\noindent ②\enspace \(60-35=25\)，\allowbreak{}\(25\times4=100\)，\allowbreak{}\(500\div100=5\)。综合算式：\rule{9em}{.4pt}。\par
}
